%%%%%
%
\section{Interpolation lin\'eaire d'une variable sur une autre base de temps.}

\itemstyle{Op\'erateur~:} \operfmt{Interpol\_irreg}
\index{Interpol\_irreg}

%%%%%OBJET%%%%
\itemstyle{Objet~:}
Cette fonction interpole une variable $X_{asyn}$, qui est dat\'ee par
le vecteur temps $Time_{asyn}$, sur une autre base de temps $Time_{reg}$ et 
traite le cas o� le nombre d'�l�ments de $X_{asyn}$ est diff�rent de celui 
de $Time_{reg}$. L'interpolation est lin\'eaire.


%%%%%ARGUMENTS%%%%
\bigskip
%%%%%ENTREES%%%%
\begin{tabarg}
\multicolumn{3}{l}{\titretab {Entr\'ees}}\\
\hline
$Time_{reg}$ & Vecteur, {\tt double} & Temps ($s$)\\
$Time_{asyn}$ & Vecteur, {\tt double} & Temps de la variable d'entr\'ee ($s$)\\
$X_{asyn}$ & Vecteur, {\tt float} & Variable \`a interpoler\\
$Size_{asyn}$ & Coeff., {\tt Int\_32} & Nombre d'�l�ments de la variable � interpoler\\
\end{tabarg}

\medskip
%%%%%SORTIES%%%%
\begin{tabarg}
\multicolumn{3}{l}{\titretab {Sorties}}\\
\hline
$X_{reg}$ & Vecteur, {\tt float} & Variable interpol\'ee \`a chaque instant de $Time_{reg}$(m\^eme unit\'e que $X_{asyn}$)\\
\end{tabarg}



%%%%%FORMULE%%%%
\itemstyle{Formule}

%%Interpol_irreg
Pour chaque valeur de $Time_{reg}$, on recherche les 2 valeurs
de $Time_{asyn}$ qui l'encadrent.
Puis on interpole lin\'eairement $X_{asyn}$ entre
les 2 instants trouv\'es.



%%%%%COMMENTAIRE%%%%
\itemstyle{Commentaire~:}
  Il faut conna�tre le nombre d'�l�ments de $X_{asyn}$ (lorsque
cet op�rateur est appel� via adc, on peut utiliser \$Arg3.NbVal)

 Cet op�rateur ne fait pas d'extrapolation, et peut donc cr�er des 
valeurs manquantes en d�but et fin de la s�rie interpoll�e, si le temps
de la premi�re valeur d'origine est apr�s le premier 
temps du r�sultat (ou sym�triquement en fin de fichier), m�me
si les �carts de temps sont tr�s petits.

%%%%%AUTEUR%%%%
\itemstyle{Auteur~:}
CNRM/GMEI/TRAMM

%%%%%REFERENCE%%%%
\itemstyle{R\'ef\'erences~:}

